% Canonical Paper I source.
This paper has established a foundational passage from sequential arithmetic
syntax to geometry.  The organizing principle is that an expression should not
be identified prematurely with its endpoint value.  An unmarked sequential tree
requires a choice of accumulator; the resulting fully marked planar tree is
equivalent to a bounded marked history.  From that history, operator evaluation
and the ACS charge shadow form separate branches, while an initial value turns an
operator into an endpoint.  The non-bijective forgetful maps lose information, and
several of the structures developed here measure precisely what is lost.

\subsection{Results of the foundational construction}

The first part of the paper classified binary expression trees with a unique
legal internal evaluation order and represented them by marked spinal
histories.  Operand-slot data made it possible to distinguish planar mirror,
temporal reversal, and path inverse.  These operations act at different levels
and need not induce the same operator comparison.

For nondegenerate bilateral one-hole contexts, projective continuation produced
M\"obius transformations.  Under the algebraic hypotheses stated there, the
elementary contexts generate \(\operatorname{PGL}_2\), while ordinary
addition and nonzero scaling form the affine, or Borel, sector stabilizing the
chosen point at infinity.  This placement is important: projective semantics
extends the operator domain, but it does not turn a pole into an admissible
ordinary arithmetic operation.

The smaller alphabet consisting of translation by $\pm\sqrt2$ and negative
inversion realizes the $q=4$ Hecke group inside this projective semantics.  Its
matrix relation $(JT)^4=1$ gives an explicit pair of different marked histories
with the same operator.  The subsequent factorization through cell stabilizers
then separates four levels that are needed downstream: history, operator, cell
coset, and endpoint.  In particular, the arithmetic words are not identified
one-to-one with tiles in a chosen Hecke cellulation.

Within the affine sector, chronological composition yielded explicit cocycle
coordinates and the continuous equation
\[
    \frac{da}{ds}
       =\mu\cos\theta+\lambda a\sin\theta.
\]
The basic upper-half-plane model realizes this equation geometrically and gives
a first regular arithmetic expression space.  Separately, the regular-value
argument shows that a zero in a regular AES lies on
a smooth codimension-one zero set.  Isolated or singular zeros must therefore
arise from a failure of those regular hypotheses, not from the regular theory
itself.  Under the additional assumptions that the metric is complete and the
manifold is connected and boundaryless, the unit-gradient rectification gives
the stronger diffeomorphism
\[
  M\cong Z(a)\times\mathbb R.
\]
Hence the zero layer is nonempty and connected.  This is not asserted for
incomplete models or manifolds with boundary, and it is not in general a
Riemannian product decomposition.

The final two chapters identified global and local forms of arithmetic order
defect.  In the accumulative commutative space, a positive add--scale history
\(\gamma\) satisfies
\[
    \nu_x(\gamma)
      =e^{M_\gamma}
        \left(x+\int_{C_\gamma}e^{-M}\,dA\right),
\]
and any charge-compatible pair of such histories has relative torsion equal to the
integral of the weighted area form over an oriented filling chain.  On the local state
space, the horizontal fields associated with
\[
    \alpha=da-\mu\,du-\lambda a\,dv
\]
satisfy
\[
    [D_u,D_v]=\mu\lambda\partial_a.
\]
The exact target-frame open defect, the source-normalized closed-loop drift, and the
infinitesimal curvature were kept distinct.  Converting the closed drift back to the
target frame gives the open defect; the two raw defects share only the stated
infinitesimal limit.

The proved logical spine may therefore be summarized as
\[
\begin{gathered}
    \text{marked planar sequential tree}
      \longleftrightarrow \text{bounded marked history}
      \longrightarrow \text{projective operator},\\
    \operatorname{PGL}_2
      \supset \operatorname{Aff}(1)
      \longrightarrow \text{affine cocycle and flow},\\
    \operatorname{Hist}_4
      \longrightarrow G_4
      \longrightarrow G_4/H_{\mathcal C}
      \longrightarrow \text{Hecke cell},\\
    \text{complete boundaryless regular AES}
      \longrightarrow Z(a)\times\mathbb R,\\
    \text{charge-compatible positive add--scale defect}
      \longrightarrow \text{ACS weighted area},\\
    \text{infinitesimal defect}
      \longrightarrow \text{contact curvature}.
\end{gathered}
\]
This is a foundation for later developments, not a claim that every arithmetic
expression, singularity, or computational phenomenon is already captured by
the affine theory.

\subsection{Interface to Paper II}

\emph{Arithmetic Expression Geometry II: Hyperbolic Real Function Theory}
begins where the present horizontal connection stops.  It develops additional
metric and analytic data: a horizontal metric,
compatible complex or almost-complex structures where appropriate,
second-order operators, boundary-value constructions, and arithmetic
holomorphicity.  None of these structures follows uniquely from the contact
form alone, and no result of that downstream function theory is used in the
proofs of Paper~I\cite{Yuan2026AEGFunctionTheory}.  The $G_4$ sublanguage
supplies an additional, explicitly arithmetic test case: an automorphic or
equivariant function used there must be
derived from the group action and must keep its analytic descent conditions
separate from ordinary arithmetic admissibility of a history at a real input.

\subsection{Interface to Paper III}

\emph{Arithmetic Expression Geometry III: Singular Zero Geometry and Tubes}
studies the boundary of the regular-zero theorem.  Its objects include singular
and multiple zeros, parameter discriminants, regular and singular zero surfaces,
and tube constructions.  Braid, monodromy, and knot structures there require
separate invariance and normalization arguments.
The Hecke interface proved here offers a source-driven alternative to imposing
a four-valent cellulation by hand: Paper~III uses a normalized Hauptmodul to
construct the corresponding singular zero dessin, then relates its sign cover
to Hecke periodic-orbit knots\cite{Yuan2026AEGZeroGeometry}.  Those analytic,
singular, and three-dimensional results remain downstream results, not results of
Paper~I.  The complete regular
splitting theorem shows that such a branching network cannot be the zero set of
one connected, complete, boundaryless regular scalar AES; any singularity,
boundary, or incompleteness used in the construction must therefore be stated
explicitly.
In particular, the existence of a tube does not by itself produce a knot
invariant, and no such general invariant is asserted here.

\subsection{Interface to Paper IV}

\emph{Arithmetic Expression Geometry IV: Projective Condensation and
Computational Complexity} returns to the full projective process--result
hierarchy.  Its scope includes bivaluation, quotient and condensation
constructions, representation complexity, and the relation between geometric
and computational resources.  The affine embedding in projective semantics
proved in this paper supplies an entry point. Paper~IV proves model-relative
quotient and residual results together with fixed-model resource calibrations, while
general machine-robust complexity claims remain open and require explicit models of
representation, input size, and computation.

\subsection{Open foundational problems}

Several questions remain at the level of foundations.
\begin{enumerate}
    \item Classify regular affine arithmetic expression spaces beyond the basic
    model on the upper half-plane, including the hypotheses under which two
    models should count as equivalent.
    \item Construct a genuinely projective replacement for the affine ACS and
    determine which part of the weighted torsion formula survives without a
    preferred point at infinity.
    \item Relate affine contact curvature to projective holonomy without
    conflating an infinitesimal bracket with a finite transport operator.
    \item Characterize the marked-history fibers of the downstream
    $G_4$ operator--zero-network realization and determine what arithmetic
    decoration, if any, survives its operator and knot quotients.
    \item Produce fully verified singular multi-zero models and identify their
    stable local normal forms and discriminants.
    \item Formulate, and then test, a quantitative relation between geometric
    history complexity, representation complexity, and computational
    complexity.
\end{enumerate}

These are research interfaces rather than consequences of the present paper.
The stable conclusion is narrower and more concrete: sequential arithmetic
histories possess projective operator semantics; their affine sector carries a
natural cocycle and flow; and the failure of additive and multiplicative order
to commute admits exact global and infinitesimal geometric descriptions.
